\documentclass[10pt,twocolumn]{article}

% --- Minimal, Overleaf-friendly preamble (no exotic packages) --------------
\usepackage[margin=0.75in]{geometry}
\usepackage[T1]{fontenc}
\usepackage{lmodern}
\usepackage{microtype}
\usepackage{graphicx}
\usepackage{booktabs}
\usepackage{amsmath}
\usepackage{xcolor}
\usepackage{caption}
\usepackage[hidelinks]{hyperref}
\captionsetup{font=small,labelfont=bf}

% Put the PNGs from docs/assets/ next to this file (or in paper/figures/).
\graphicspath{{figures/}{../docs/assets/}}

\newcommand{\pc}{\%}  % readable percent in text

\title{\vspace{-2.5em}\textbf{Generalization and Evaluation Rigor in Deep\\
Reinforcement Learning: A From-Scratch Study on a\\
Financial-Market Testbed}}
\author{Daniel Lichtenberger\\
\small\texttt{github.com/Danny-397/RL-for-Crypto-and-stocks-}}
\date{\vspace{-1.2em}\small 2026}

\begin{document}
\maketitle

\begin{abstract}
\noindent Deep reinforcement learning is unusually easy to fool yourself with:
agents memorize their training trajectory, and a single lucky random seed can
masquerade as a real result. We study both failure modes from scratch --- a
PyTorch implementation of Proximal Policy Optimization (PPO), two custom
Gymnasium environments, and a full evaluation suite --- using financial markets
as a deliberately hard testbed: non-stationary, partially observable, with a
noisy reward and near--random-walk signal. We ask (RQ1) whether one fixed recipe
generalizes across two market regimes; (RQ2) how severely an agent overfits a
single price trajectory and whether \emph{domain randomization} repairs it;
(RQ3) whether an apparent out-of-sample edge survives multi-seed significance
testing; and (RQ4) whether the same holds for cross-sectional allocation. Our
headline result is negative, and that is the point: on real markets the agent
has \emph{no seed-robust edge} over buy-and-hold, and our own significance
tooling catches it: an earlier build of this study published a $+275\pc$
single-seed crypto run that survived neither reseeding nor a two-month extension
of the evaluation window. We further
introduce a \textbf{surrogate-data falsification test} that separates ``the agent
is weak'' from ``there is no signal to find,'' and validate it with a positive
control on synthetic data with known structure. In doing so the project
reproduces, in a new domain, the central methodological finding of
Henderson~et~al.\ (2018) --- that single-run RL evaluation is unreliable. We then
apply that finding to our own ablation, whose domain-randomized arm is unseeded by
construction and so cannot be reproduced from a single run: we report both arms as
distributions over five seeds, where all four confidence intervals exclude zero.
\end{abstract}

\section{Introduction}
Reinforcement learning (RL) research has a reproducibility problem: reported
performance swings wildly across random seeds, implementation details, and
evaluation choices \cite{henderson2018}. Trading is a domain where these hazards
are especially acute --- the data-generating process drifts (non-stationarity),
price alone hides the true state (partial observability), the reward is noisy,
and a single historical path is trivial to memorize. That makes markets an
excellent \emph{stress test} for the general RL questions of generalization and
honest evaluation, rather than a get-rich scheme.

We build the entire stack from the algorithm up --- PPO \cite{ppo} with
Generalized Advantage Estimation \cite{gae}, two market environments, a
28-feature pipeline, and the evaluation harness --- so that every experiment can
be instrumented and trusted. Our contributions are:
\begin{itemize}\itemsep2pt
  \item A controlled study showing an RL agent memorizes a single price path
  catastrophically, and that \textbf{domain randomization}
  \cite{tobin2017,cobbe2019} is the single change that flips held-out return from
  reliably negative to reliably positive --- across five seeds all four confidence
  intervals exclude zero, in both markets (\S\ref{sec:ablation}).
  \item A \textbf{multi-seed significance protocol} (bootstrap CIs + paired
  permutation tests) that reproduces the Henderson~et~al.\ \cite{henderson2018} finding in a new
  domain: a favorable single-seed ``win'' does not survive resampling
  (\S\ref{sec:sig}).
  \item An original \textbf{surrogate-data falsification test} that distinguishes
  agent weakness from signal absence, with a synthetic positive control that
  confirms the test has statistical power (\S\ref{sec:surrogate}).
\end{itemize}

\section{Method}
\paragraph{Algorithm.} A single \texttt{PPOAgent} (shared-trunk actor--critic)
implements the clipped surrogate objective, GAE advantages, entropy
regularization, orthogonal initialization, gradient clipping, and a Welford
\cite{welford1962} running observation normalizer that is exported and reused at
serving time. A recurrent (LSTM) actor--critic trains through the same machinery
with truncated backpropagation through time.

\paragraph{Environments.} \texttt{StockTradingEnv} and \texttt{CryptoTradingEnv}
subclass one \texttt{BaseTradingEnv} so accounting, transaction costs, slippage,
and reward can never silently diverge. The action is a continuous target position
in $[-1,1]$ (fraction of equity; negative $=$ short). The observation is a rolling
window of 28 stationary features spanning multi-horizon momentum, moving-average
ratios, RSI/MACD, Bollinger \%B, Donchian position, ATR and volatility-regime
signals, drawdown-from-high, volume microstructure, and cross-asset context
(relative strength vs.\ SPY/BTC and the index's own trend).

\paragraph{Reward.} Two selectable forms: a risk-aware return net of transaction
costs with drawdown and turnover penalties, or the Differential Sharpe Ratio
\cite{moody2001}, an online per-step approximation to the change in the Sharpe
ratio.

\paragraph{Evaluation protocol.} Data are split chronologically; feature scalers
are fit on the training slice only. Agents are scored on an untouched test split
by total return, annualized Sharpe, and maximum drawdown, against buy-and-hold,
random, and moving-average-crossover baselines run through the identical cost
model. Every headline claim is repeated across independent seeds and reported as a
bootstrap 95\pc\ confidence interval, with a paired permutation test against
buy-and-hold.

\section{Experiments}

\subsection{RQ1: Regime transfer}
Holding the architecture and hyperparameters fixed, we train one agent per market
using only market-specific environment dynamics (cost, slippage, volatility). The
same recipe produces coherent, market-appropriate behavior on both equities and
crypto, confirming the unified agent is genuinely market-agnostic and that
downstream differences are attributable to the environments, not the code.

\subsection{RQ2: Overfitting and domain randomization}
\label{sec:ablation}
Training on a \emph{single} price series is a trap: the agent memorizes it. We
train two otherwise-identical agents --- one on a fixed synthetic path, one drawing
a fresh path every episode (domain randomization) --- and measure in-sample vs.\
out-of-sample (30 unseen paths) return (Table~\ref{tab:ablation}).

\begin{table}[t]\centering\footnotesize
\caption{Domain-randomization ablation, 5 seeds $\times$ 60k steps; all cells are
percent returns. Held-out is the mean over 30 unseen paths --- identical for every
arm and seed --- with a bootstrap 95\pc\ CI across seeds. In-sample is a
\emph{range} over seeds, not a mean: it is one leveraged compounding return on a
memorized path and spans orders of magnitude. The domain-randomized arm's in-sample
figure is its return on the very path the single-path agent memorized.}
\label{tab:ablation}
\setlength{\tabcolsep}{4pt}
\begin{tabular}{@{}llrr@{}}
\toprule
Market & Train & In-sample & Held-out [95\pc\ CI]\\
\midrule
% BEGIN GENERATED: tex-ablation
Stock  & single & $+17k$--$+30k$ & $-26$ {\scriptsize$[-37,-9]$}\\
Stock  & domain & $-1$--$+128$ & $\mathbf{+47}$ {\scriptsize$[+31,+64]$}\\
Crypto & single & $+322k$--$+4.0M$ & $-51$ {\scriptsize$[-62,-39]$}\\
Crypto & domain & $+5$--$+154$ & $\mathbf{+129}$ {\scriptsize$[+82,+185]$}\\
% END GENERATED: tex-ablation
\bottomrule
\end{tabular}
\end{table}

\begin{figure}[t]\centering
\includegraphics[width=\columnwidth]{fig_ablation}
\caption{Held-out return by training regime. Bars are the mean over seeds, whiskers
the bootstrap 95\pc\ CI, and dots the individual seeds --- so the reader sees $n$
and the spread rather than a bar alone. The in-sample range is annotated rather
than plotted: it spans several orders of magnitude and would flatten everything
else to a line.}
\label{fig:ablation}
\end{figure}

The astronomical in-sample figures are not a win --- they are the overfitting
signature. Domain randomization is the single change that turns a memorizer into a
generalizer, echoing sim-to-real transfer \cite{tobin2017} and RL generalization
benchmarks \cite{cobbe2019}.

Across five seeds all four intervals exclude zero: single-path training is
reliably unprofitable out-of-sample and domain randomization is reliably
profitable, in both markets. The effect is therefore robust to seed choice, which
is the property the rest of this paper spends its time showing that real-market
edges lack.

Two caveats the distribution makes visible and a single run would hide. Equities
are the weaker case --- one of five single-path seeds returned $+8\pc$ rather than
a loss, and the interval reaches to $-9\pc$, so ``reliably negative'' there rests
on a margin, not a chasm. And the crypto domain-randomized mean is inflated by one
seed at $+307\pc$ against a cluster near $+60\pc$; the interval is correspondingly
wide. Figure~\ref{fig:ablation} plots the individual seeds for exactly this reason.

\subsection{RQ3: Does the edge survive multiple seeds?}
\label{sec:sig}
A single backtest is an anecdote. On synthetic markets where a signal provably
exists, the agent is repeatably profitable across 5 seeds --- its bootstrap
interval excludes zero in both markets --- yet remains statistically
indistinguishable from buy-and-hold under a paired permutation test. On
\emph{real} markets the two markets fail differently (Table~\ref{tab:real}). On
equities the agent is significantly \emph{worse} than the mega-cap bull. On crypto
the mean across seeds is positive and its interval excludes zero, so the agent did
make money on average; what it did not do is beat buy-and-hold, which the paired
test cannot distinguish it from. The per-seed returns span more than an order of
magnitude, and with 5 seeds a sign-flip test on that axis cannot reach $p{\le}0.05$
even in principle --- its floor is $2/2^5{=}0.0625$. An earlier build of this study
(commit \texttt{d4c0ef9}) published the crypto agent at $+275\pc$ against
buy-and-hold's $+19\pc$; rebuilding the identical recipe --- same code, same seed,
same budget --- against a data snapshot pinned two months later erased that
headline entirely. The run was not merely seed-lucky but window-lucky.

\begin{table}[t]\centering\small
\caption{Real-market multi-seed significance (5 seeds, held-out walk-forward). The
lucky single-seed crypto ``win'' does not survive resampling.}
\label{tab:real}
\begin{tabular}{@{}lrrl@{}}
\toprule
Market & Agent (95\pc\ CI) & vs.\ B\&H & Verdict\\
\midrule
% BEGIN GENERATED: tex-significance
Crypto & $+79.7\pc$ $[+1,+163]$ & $+46\pc$ & indist. ($p{=}0.8191$)\\
Stock  & $-21.5\pc$ $[-29,-15]$ & $-261\pc$ & worse ($p{=}0.0021$)\\
% END GENERATED: tex-significance
\bottomrule
\end{tabular}
\end{table}

This is exactly the warning of Henderson~et~al.\ \cite{henderson2018}: a single run is not evidence.
A naive project ships the lucky table; the significance test is what catches it.
\emph{Reporting the catch is the contribution.}

\subsection{RQ (original): A surrogate-data falsification test}
\label{sec:surrogate}
The results above show the agent does not beat the market --- but is that because
the agent is weak, or because there is no exploitable structure to find? We adapt
\textbf{surrogate-data testing} from nonlinear time-series analysis
\cite{theiler1992}. A \emph{return-shuffled surrogate} randomly permutes a series'
daily log-returns and re-integrates them. Because a sum is permutation-invariant,
the surrogate ends at the identical price --- \emph{buy-and-hold is unchanged} ---
but momentum, autocorrelation, and volatility clustering are destroyed. We then
train and evaluate the same recipe on structured data and on its surrogate and
compare the agent's edge over buy-and-hold.

The logic is a controlled falsification. On synthetic data with a \emph{known}
AR(1) momentum signal, a valid test must find $\text{edge}_{\text{structured}} \gg
\text{edge}_{\text{surrogate}}$: the agent should exploit the signal and lose that
edge once shuffling removes it. This is our positive control
(Table~\ref{tab:surrogate}), and it confirms the test has power. Applied to real
markets, the test finds \emph{no} significant difference between real and
return-shuffled data: on equities the agent does no better on real prices than on
structure-free surrogates ($\Delta$ nowhere near significant), the clean
null predicted by \S\ref{sec:sig}. The crypto arm is inconclusive --- heavy-tailed
returns produce occasional blow-ups on reshuffled paths that inflate the surrogate
loss and dominate the mean, and the gap is not significant; a median-based or
exposure-capped
statistic is the natural remedy. Neither market shows demonstrable exploited
structure, corroborating \S\ref{sec:sig} from a fresh angle and turning ``markets
are efficient'' from an appeal to authority into a falsifiable test.

\begin{table}[t]\centering\small
\caption{Surrogate-data test, synthetic positive control (edge $=$ agent $-$
buy-and-hold return, mean over 10 held-out paths, 5 seeds, 60k steps; $p$ from a
paired permutation test of structured vs.\ surrogate edge). When a known signal
exists the agent earns a positive edge; destroying it by shuffling collapses the
edge sharply negative --- the test has power in both markets.}
\label{tab:surrogate}
\begin{tabular}{@{}lrrrr@{}}
\toprule
Market & Structured & Surrogate & $\Delta$ & $p$\\
\midrule
% BEGIN GENERATED: tex-surrogate-control
Crypto & $+68.6\pc$ & $-55.7\pc$ & $+124.3\pc$ & $0.0075$\\
Stock  & $+1.8\pc$ & $-47.7\pc$ & $+49.5\pc$ & $0.0026$\\
% END GENERATED: tex-surrogate-control
\bottomrule
\end{tabular}
\end{table}

\subsection{RQ4: Cross-sectional allocation}
Single-asset timing is only half the problem. The same PPO agent generalizes to a
\texttt{PortfolioTradingEnv} that observes a whole basket and emits an
$N$-dimensional long/short weight vector under a gross-exposure budget, benchmarked
against an equal-weight ($1/N$) portfolio and a cross-sectional momentum factor. On
equities the learned allocator clears both random weights and the momentum factor
yet is beaten by equal-weight by more than 150 points of return; on crypto it is
level with random weights on return, ahead only on a risk-adjusted basis. The
contribution is the working, budget-constrained RL allocator and its
apples-to-apples evaluation --- the same lesson at a harder problem: naive
diversification is hard to beat when the per-asset signal is weak.

\section{Related Work}
This project reproduces and cross-domain stress-tests results the field already
considers important. Henderson~et~al.\ \cite{henderson2018} document the unreliability of
single-run deep-RL evaluation; \S\ref{sec:sig} reproduces it on markets.
Tobin~et~al.\ \cite{tobin2017} introduce domain randomization for sim-to-real transfer and
Cobbe~et~al.\ \cite{cobbe2019} formalize train/test generalization in RL; \S\ref{sec:ablation}
measures the gap domain randomization closes. The algorithm is PPO \cite{ppo} with
GAE \cite{gae}; the Differential Sharpe Ratio reward is due to
Moody and Saffell \cite{moody2001}. The negative real-market result is consistent with weak-form
market efficiency \cite{fama1970}, and surrogate-data testing follows
Theiler~et~al.\ \cite{theiler1992}.

\section{Limitations and Conclusion}
The bottleneck is signal, not the agent: raw daily OHLCV carries little
exploitable structure, so richer exogenous data (macro series, sentiment) --- not a
bigger network --- is the natural next lever. Multi-fold walk-forward and a
head-to-head feed-forward-vs-LSTM study are supported by the codebase and left as
future work.

A second limitation concerns the ablation of \S\ref{sec:ablation}, and we state it
plainly because we found it in our own results. The two arms are not equally
reproducible, and the asymmetry is structural rather than incidental. The
single-path arm draws its training series with an explicit seed, and repeated
invocations reproduce it exactly. The domain-randomized arm, by construction,
draws a \emph{fresh unseeded} series every episode --- that is what domain
randomization means --- so the run seed does not determine its data stream and no
single run of it is reproducible even in principle. Any comparison between the
arms that quotes one run of each is therefore not a controlled comparison.

This is why \S\ref{sec:ablation} reports each arm as a distribution over seeds
with bootstrap intervals, and why its conclusion is stated in terms of held-out
return --- a mean over 30 fixed evaluation paths --- rather than the single-path
in-sample figure, which averages nothing and compounds along one memorized
trajectory. Training additionally does not run under
\texttt{torch.use\_deterministic\_algorithms}, so exact bitwise reproducibility is
not guaranteed for either arm across machines. An earlier version of this table
quoted a single run per cell; that version was also stale with respect to a change
in the observation feature set, which is the ordinary way such numbers rot and an
argument for regenerating them from a script rather than transcribing them.

Our conclusion is deliberately unglamorous and, we argue, more
valuable than a suspicious backtest: built and evaluated honestly, a from-scratch
deep-RL trader finds no reliable edge on real markets --- and the methodology that
establishes this rigorously is the actual result.

\begin{thebibliography}{9}\small
\bibitem{henderson2018} P.~Henderson, R.~Islam, P.~Bachman, J.~Pineau, D.~Precup,
  D.~Meger. Deep Reinforcement Learning that Matters. \emph{AAAI}, 2018.
\bibitem{ppo} J.~Schulman, F.~Wolski, P.~Dhariwal, A.~Radford, O.~Klimov.
  Proximal Policy Optimization Algorithms. \emph{arXiv:1707.06347}, 2017.
\bibitem{gae} J.~Schulman, P.~Moritz, S.~Levine, M.~Jordan, P.~Abbeel.
  High-Dimensional Continuous Control Using GAE. \emph{ICLR}, 2016.
\bibitem{tobin2017} J.~Tobin et al. Domain Randomization for Transferring Deep
  Neural Networks from Simulation to the Real World. \emph{IROS}, 2017.
\bibitem{cobbe2019} K.~Cobbe, O.~Klimov, C.~Hesse, T.~Kim, J.~Schulman.
  Quantifying Generalization in Reinforcement Learning. \emph{ICML}, 2019.
\bibitem{moody2001} J.~Moody, M.~Saffell. Learning to Trade via Direct
  Reinforcement. \emph{IEEE Trans.\ Neural Networks}, 2001.
\bibitem{theiler1992} J.~Theiler et al. Testing for Nonlinearity in Time Series:
  The Method of Surrogate Data. \emph{Physica D}, 1992.
\bibitem{fama1970} E.~Fama. Efficient Capital Markets. \emph{J.\ Finance}, 1970.
\bibitem{welford1962} B.~P.~Welford. Note on a Method for Calculating Corrected
  Sums of Squares and Products. \emph{Technometrics}, 1962.
\end{thebibliography}

\end{document}
